\section{H1预告是分母更新，而非自动上调倍数}
7月14日公告把研究分母从Q1观察推进到H1区间：归母净利润CNY1.70--2.40亿元、扣非净利润CNY2.30--3.30亿元，收入同比增长56--66\%。公告同时说明，云创数据股权公允价值变动对本期税前利润影响约为负CNY0.91亿元。归母和扣非差距很大，因此本报告以扣非利润作为经营能力分母，以归母利润作为股东收益压力测试，不把一次性公允价值变化当成核心经营趋势。

\begin{exhibitbox}[表4-1\quad 2026E盈利桥]
\begin{tabularx}{\exhibitboxwidth}{L{2.8cm}R{1.8cm}R{1.8cm}R{1.8cm}X}
\textbf{指标CNY100m} & \textbf{熊} & \textbf{基准} & \textbf{牛} & \textbf{桥接方式} \\
H1扣非净利 & 2.30 & 2.80 & 3.30 & 公司官方预告 \\
H2扣非净利 & 1.10 & 2.40 & 3.70 & AStock情景假设，非公司指引 \\
2026E扣非净利 & 3.40 & 5.20 & 7.00 & H1+H2 \\
2026E扣非EPS & 0.276 & 0.423 & 0.569 & 扣非净利/12.299945亿股 \\
2026E收入 & 47.49 & 49.01 & 50.54 & H1收入中点附近的情景外推 \\
\end{tabularx}
\sourcenote{data/growth\_driver\_model.json；data/current\_valuation\_model\_20260714.json}
\end{exhibitbox}

\section{增长驱动与未完成字段}
模型把“光纤需求量价齐升”视为已披露的周期恢复驱动，把安全子公司并表和经营改善视为已披露的利润贡献，把AI/数据中心曝光视为期权信用。因为H1没有产品级收入、ASP、利用率、客户分配、订单金额和分部毛利，不能用“单位×ASP×收入确认率”伪造精确的光纤EPS。增长模型的门禁状态为CONDITIONAL，而不是PASS。

\section{下一季验证阈值}
\begin{enumerate}
\item 根据H1区间反推，Q2扣非净利润约CNY1.27--2.27亿元；基准情景要求接近CNY2.40亿元。
\item 正式半年报应解释光纤收入、毛利率、存货和订单转换，不能只有公司层净利润。
\item 经营现金流需要随利润改善，若持续为负，估值倍数不能继续扩张。
\item 2025年合并毛利率为19.45\%，若H1明显低于此水平，必须有产品结构解释。
\end{enumerate}
